% Part: many-valued-logic
% Chapter: syntax-and-semantics
% Section: matrices

\documentclass[../../../include/open-logic-section]{subfiles}

\begin{document}

\olfileid[pt]{mvl}{syn}{mat}

\olsection{Matrizes}

Uma lógica multivalorada é definida por sua linguagem, seu conjunto de valores de
verdade~$V$, um subconjunto de valores de verdade designados e funções de verdade
para seus conectivos. Em conjunto, esses elementos são chamados de \emph{matriz}.  

\begin{defn}[Matriz]
\ollabel{defn:matrix}
Uma \emph{matriz} para a lógica~$\Log L$ consiste em:
\begin{enumerate}
\item um conjunto de conectivos que constitui uma linguagem~$\Lang L$;
\item um conjunto $V \neq \emptyset$ de valores de verdade;
\item um conjunto $V^+ \subseteq V$ de valores de verdade designados;
\item para cada conectivo $n$-ário $\star$ de $\Lang L$, uma função de
verdade~$\tf{\star} : V^n \to V$. Se $n = 0$, então $\tf{\star}$ é
simplesmente um elemento de~$V$.
\end{enumerate}
\end{defn}

\begin{ex}
A matriz da lógica clássica~\LogCL{} consiste em:
\begin{enumerate}
  \item A linguagem proposicional padrão $\Lang L_0$, com
  $\lfalse$, $\lnot$, $\land$, $\lor$, $\lif$.
  \item O conjunto de valores de verdade $V = \{\True, \False\}$.
  \item $\True$ é o único valor designado, isto é, $V^+ = \{\True\}$.
  \item Para $\lfalse$, temos $\tf{\lfalse} = \False$. As demais
  funções de verdade são dadas pelas tabelas-verdade usuais (veja
  \olref{fig:tf-CL}).
\end{enumerate}
\begin{figure}
  \begin{center}
      \begin{tabular}{c|c} 
        $\tf{\lnot}$ & \\ 
        \hline  
        $\True$ & $\False$ \\ 
        $\False$ & $\True$ 
      \end{tabular}
      \quad
      \begin{tabular}{c|cc} 
        $\tf{\land}$ & $\True$ & $\False$ \\ 
        \hline 
        $\True$ & $\True$ & $\False$ \\ 
        $\False$ & $\False$ & $\False$ 
      \end{tabular}
      \quad
      \begin{tabular}{c|cc} 
        $\tf{\lor}$ & $\True$ & $\False$ \\ 
        \hline 
        $\True$ & $\True$ & $\True$ \\ 
        $\False$ & $\True$ & $\False$ 
      \end{tabular}
      \quad
      \begin{tabular}{c|cc} 
        $\tf{\lif}$ & $\True$ & $\False$ \\ 
        \hline 
        $\True$ & $\True$ & $\False$ \\ 
        $\False$ & $\True$ & $\True$ 
      \end{tabular}
    \end{center} 
    \caption{Funções de verdade da lógica clássica~$\LogCL$.}
    \ollabel{fig:tf-CL}
  \end{figure}
  \end{ex}

\end{document}
